\documentclass[../main.tex]{subfiles}

\begin{document}

\chapter*{Introduction générale}
\addcontentsline{toc}{chapter}{Introduction générale}

Le continent africain représente aujourd'hui l'une des dernières grandes frontières
de croissance pour le secteur de la réassurance mondiale. Avec un taux de pénétration
de l'assurance inférieur à 3\,\% du PIB en Afrique subsaharienne --- contre plus de
7\,\% au niveau mondial \citep{SwissRe2023} --- et des économies portées par une
croissance démographique soutenue, des projets d'infrastructure massifs et une
exposition croissante aux risques climatiques, l'asymétrie entre besoins de couverture
et capacité assurantielle déployée constitue pour un réassureur une opportunité
commerciale et technique de premier ordre.

C'est dans cette logique qu'Atlantic Re --- réassureur international, filiale du
Groupe CDG opérant dans plus de 60 pays auprès de plus de 400 clients --- a engagé,
dans le cadre de son plan stratégique \textbf{Reach2030}, une stratégie d'expansion
ciblée sur le continent africain. Avec un chiffre d'affaires atteignant
\textbf{4~039 MDH} et un résultat net de \textbf{419 MDH} au titre de l'exercice
2025, un ratio combiné de \textbf{92,3\,\%} et un ratio de solvabilité S2 de
\textbf{215\,\%}, Atlantic Re dispose des bases financières et techniques pour
accélérer cette expansion. L'enjeu est désormais concret : identifier, parmi les
marchés africains Vie et Non-Vie aux profils très hétérogènes, ceux qui offrent la
meilleure combinaison de potentiel de primes, de rentabilité technique, de stabilité
réglementaire et de dynamisme économique.

Ce projet de fin d'études s'inscrit dans cette commande stratégique. Il est toutefois
important d'en comprendre la trajectoire réelle, car elle reflète une démarche
d'ingénieur autant que de modélisateur. Débuté comme un projet de modélisation pure,
il a évolué, au contact du terrain, vers un projet de \textbf{modélisation et de
digitalisation} : c'est lors de l'immersion métier dans le département Business
Développement que nous avons détecté un problème structurel non anticipé --- l'absence
de tout outil permettant aux équipes d'exploiter la donnée interne --- et que nous
avons proposé, de notre propre initiative, d'y apporter une réponse digitale en
complément du modèle externe. Cette décision proactive constitue l'un des apports
distinctifs de ce stage.

Le rapport s'organise en quatre chapitres. Le \textbf{Chapitre~1} présente Atlantic
Re, son contexte stratégique africain et l'objectif initial du PFE. Le
\textbf{Chapitre~2} expose les fondements métier de la réassurance, l'analyse des
besoins --- incluant les besoins identifiés lors de l'immersion et qui ont conduit
à l'élargissement du périmètre --- et la conception des deux axes selon une approche
itérative. Le \textbf{Chapitre~3} présente et justifie l'ensemble des outils,
algorithmes et choix d'architecture technique. Le \textbf{Chapitre~4} détaille les
réalisations, les résultats obtenus et leur validation.

\clearpage

% ============================================================
%  CHAPITRE 1 : CONTEXTE GÉNÉRAL DU PROJET
% ============================================================

\end{document}
